% Options for packages loaded elsewhere
\PassOptionsToPackage{unicode}{hyperref}
\PassOptionsToPackage{hyphens}{url}
\documentclass[
]{article}
\usepackage{xcolor}
\usepackage{amsmath,amssymb}
\setcounter{secnumdepth}{5}
\usepackage{iftex}
\ifPDFTeX
  \usepackage[T1]{fontenc}
  \usepackage[utf8]{inputenc}
  \usepackage{textcomp} % provide euro and other symbols
\else % if luatex or xetex
  \usepackage{unicode-math} % this also loads fontspec
  \defaultfontfeatures{Scale=MatchLowercase}
  \defaultfontfeatures[\rmfamily]{Ligatures=TeX,Scale=1}
\fi
\usepackage{lmodern}
\ifPDFTeX\else
  % xetex/luatex font selection
\fi
% Use upquote if available, for straight quotes in verbatim environments
\IfFileExists{upquote.sty}{\usepackage{upquote}}{}
\IfFileExists{microtype.sty}{% use microtype if available
  \usepackage[]{microtype}
  \UseMicrotypeSet[protrusion]{basicmath} % disable protrusion for tt fonts
}{}
\makeatletter
\@ifundefined{KOMAClassName}{% if non-KOMA class
  \IfFileExists{parskip.sty}{%
    \usepackage{parskip}
  }{% else
    \setlength{\parindent}{0pt}
    \setlength{\parskip}{6pt plus 2pt minus 1pt}}
}{% if KOMA class
  \KOMAoptions{parskip=half}}
\makeatother
\usepackage{color}
\usepackage{fancyvrb}
\newcommand{\VerbBar}{|}
\newcommand{\VERB}{\Verb[commandchars=\\\{\}]}
\DefineVerbatimEnvironment{Highlighting}{Verbatim}{commandchars=\\\{\}}
% Add ',fontsize=\small' for more characters per line
\newenvironment{Shaded}{}{}
\newcommand{\AlertTok}[1]{\textcolor[rgb]{1.00,0.00,0.00}{\textbf{#1}}}
\newcommand{\AnnotationTok}[1]{\textcolor[rgb]{0.38,0.63,0.69}{\textbf{\textit{#1}}}}
\newcommand{\AttributeTok}[1]{\textcolor[rgb]{0.49,0.56,0.16}{#1}}
\newcommand{\BaseNTok}[1]{\textcolor[rgb]{0.25,0.63,0.44}{#1}}
\newcommand{\BuiltInTok}[1]{\textcolor[rgb]{0.00,0.50,0.00}{#1}}
\newcommand{\CharTok}[1]{\textcolor[rgb]{0.25,0.44,0.63}{#1}}
\newcommand{\CommentTok}[1]{\textcolor[rgb]{0.38,0.63,0.69}{\textit{#1}}}
\newcommand{\CommentVarTok}[1]{\textcolor[rgb]{0.38,0.63,0.69}{\textbf{\textit{#1}}}}
\newcommand{\ConstantTok}[1]{\textcolor[rgb]{0.53,0.00,0.00}{#1}}
\newcommand{\ControlFlowTok}[1]{\textcolor[rgb]{0.00,0.44,0.13}{\textbf{#1}}}
\newcommand{\DataTypeTok}[1]{\textcolor[rgb]{0.56,0.13,0.00}{#1}}
\newcommand{\DecValTok}[1]{\textcolor[rgb]{0.25,0.63,0.44}{#1}}
\newcommand{\DocumentationTok}[1]{\textcolor[rgb]{0.73,0.13,0.13}{\textit{#1}}}
\newcommand{\ErrorTok}[1]{\textcolor[rgb]{1.00,0.00,0.00}{\textbf{#1}}}
\newcommand{\ExtensionTok}[1]{#1}
\newcommand{\FloatTok}[1]{\textcolor[rgb]{0.25,0.63,0.44}{#1}}
\newcommand{\FunctionTok}[1]{\textcolor[rgb]{0.02,0.16,0.49}{#1}}
\newcommand{\ImportTok}[1]{\textcolor[rgb]{0.00,0.50,0.00}{\textbf{#1}}}
\newcommand{\InformationTok}[1]{\textcolor[rgb]{0.38,0.63,0.69}{\textbf{\textit{#1}}}}
\newcommand{\KeywordTok}[1]{\textcolor[rgb]{0.00,0.44,0.13}{\textbf{#1}}}
\newcommand{\NormalTok}[1]{#1}
\newcommand{\OperatorTok}[1]{\textcolor[rgb]{0.40,0.40,0.40}{#1}}
\newcommand{\OtherTok}[1]{\textcolor[rgb]{0.00,0.44,0.13}{#1}}
\newcommand{\PreprocessorTok}[1]{\textcolor[rgb]{0.74,0.48,0.00}{#1}}
\newcommand{\RegionMarkerTok}[1]{#1}
\newcommand{\SpecialCharTok}[1]{\textcolor[rgb]{0.25,0.44,0.63}{#1}}
\newcommand{\SpecialStringTok}[1]{\textcolor[rgb]{0.73,0.40,0.53}{#1}}
\newcommand{\StringTok}[1]{\textcolor[rgb]{0.25,0.44,0.63}{#1}}
\newcommand{\VariableTok}[1]{\textcolor[rgb]{0.10,0.09,0.49}{#1}}
\newcommand{\VerbatimStringTok}[1]{\textcolor[rgb]{0.25,0.44,0.63}{#1}}
\newcommand{\WarningTok}[1]{\textcolor[rgb]{0.38,0.63,0.69}{\textbf{\textit{#1}}}}
\setlength{\emergencystretch}{3em} % prevent overfull lines
\providecommand{\tightlist}{%
  \setlength{\itemsep}{0pt}\setlength{\parskip}{0pt}}
\usepackage{bookmark}
\IfFileExists{xurl.sty}{\usepackage{xurl}}{} % add URL line breaks if available
\urlstyle{same}
\hypersetup{
  hidelinks,
  pdfcreator={LaTeX via pandoc}}

\author{}
\date{}

\begin{document}

{
\setcounter{tocdepth}{3}
\tableofcontents
}
\section{Reading note: Section 2.1 - Non-Abelian anyons from
gauging}\label{reading-note-section-2.1---non-abelian-anyons-from-gauging}

\textbf{Paper:} \texttt{arXiv-2602.12228v1}, Section 2.1, ``Non-Abelian
Anyons from Gauging''\\
\textbf{Aim:} explain the \emph{physical claim} of this subsection
before the lattice derivation in Section 2.2.

\subsection{The one-sentence story}\label{the-one-sentence-story}

Start with two ordinary (Abelian) toric codes. A global operation
changes the \emph{type} of certain anyons. When we \textbf{gauge} that
operation, its boundaries become physical, mobile gauge-flux strings.
Taking an anyon around one of those fluxes changes its type rather than
merely multiplying its wavefunction by a phase. That is the mechanism
that produces non-Abelian braiding.

\subsection{What has already been set
up}\label{what-has-already-been-set-up}

The beginning of Section 2 places two square-lattice toric codes on top
of one another, called \textbf{B} (blue) and \textbf{C} (orange). Each
has two basic anyon types:

\begin{itemize}
\tightlist
\item
  \texttt{e} (``electric charge''): a violation of an
  \texttt{X}-type/star stabilizer;
\item
  \texttt{m} (``magnetic flux''): a violation of a
  \texttt{Z}-type/plaquette stabilizer.
\end{itemize}

\texttt{m\_B\ e\_C} means an \texttt{m} anyon in copy B and an
\texttt{e} anyon in copy C at the same location: a composite excitation.
Before gauging, all anyons of this two-copy toric code are Abelian:
fusing two specified types has one definite outcome.

The authors define a physical finite-depth circuit \texttt{U\_CZ}, made
from controlled-Z gates between nearby B and C qubits. On the
ground-state (encoded-qubit) subspace it acts as a logical controlled-Z.
Its essential action is

\[
\bar X_B \mapsto \bar X_B\bar Z_C,\qquad
\bar X_C \mapsto \bar Z_B\bar X_C,\qquad
\bar Z_B\mapsto\bar Z_B,\quad \bar Z_C\mapsto\bar Z_C.
\]

This is the familiar controlled-Z identity
\texttt{X\ -\textgreater{}\ XZ} on the other qubit, now applied to
logical qubits. Since controlled-Z squared is the identity, this is a
\texttt{Z\_2} symmetry.

\subsubsection{Why it is called an ``emergent
symmetry''}\label{why-it-is-called-an-emergent-symmetry}

\texttt{U\_CZ} does not leave every individual toric-code Hamiltonian
term unchanged. For example, it turns an \texttt{X} stabilizer in B into
that stabilizer times a nearby \texttt{Z} stabilizer in C. But
stabilizers have value \texttt{+1} in the ground space, so the circuit
maps the ground space to itself. It is a symmetry of the encoded
low-energy information, rather than a literal symmetry of every
microscopic term.

\subsection{Step 1: the symmetry changes anyon
type}\label{step-1-the-symmetry-changes-anyon-type}

A magnetic anyon is the endpoint of an \texttt{X}-string. When such a
string enters a region on which controlled-Z was applied, its \texttt{X}
part picks up a \texttt{Z} part in the other copy. The endpoint of that
new \texttt{Z} string is an electric anyon. Consequently,

\[
g:\quad m_B\longmapsto m_B e_C,\qquad m_C\longmapsto m_C e_B,
\]

while \texttt{e\_B} and \texttt{e\_C} are unchanged. The paper calls
\texttt{g} an \textbf{anyon automorphism}: it relabels allowed anyon
types while preserving their fusion rules. Applying it twice returns
every type, because \texttt{e\_B\^{}2=e\_C\^{}2=1}.

This is the crucial fact. The symmetry does not merely attach a phase to
\texttt{m\_B}; it turns \texttt{m\_B} into a physically distinct type,
\texttt{m\_B\ e\_C}.

\subsection{Step 2: a symmetry patch is a type-changing
boundary}\label{step-2-a-symmetry-patch-is-a-type-changing-boundary}

Imagine applying \texttt{U\_CZ} only inside a finite region. This is a
\textbf{symmetry patch}. At its boundary, the inside and outside anyon
labels are related by \texttt{g}. Thus an anyon crossing the boundary is
relabelled:

\begin{Shaded}
\begin{Highlighting}[]
\NormalTok{outside: m\_B       {-}{-}cross the patch boundary{-}{-}\textgreater{}       inside: m\_B e\_C}
\end{Highlighting}
\end{Shaded}

This is what \texttt{patch.pdf} depicts. The boundary is a
\textbf{symmetry-defect line}. At this point it is externally imposed:
we chose where to apply the patch.

\subsection{Step 3: gauging makes that boundary
dynamical}\label{step-3-gauging-makes-that-boundary-dynamical}

To \textbf{gauge} the global \texttt{Z\_2} symmetry, the authors add
qubits on a red \textbf{A} sublattice which record whether the symmetry
has been applied locally, then impose a local gauge constraint. A useful
slogan is:

\begin{quote}
Before gauging, a symmetry patch is something we deliberately apply.
After gauging, its boundary is a physical gauge-field configuration that
the system can contain.
\end{quote}

In the gauged model, a closed A-gauge-field loop carries the same
information as the boundary of a controlled-Z patch. Crossing it still
performs \texttt{g} on a B/C anyon. This is why
\texttt{gauge\_field.pdf} is reinterpreted as the symmetry-defect loop
in \texttt{symmetry\_defect.pdf}.

Such a loop can be opened into a string. Its endpoints are new
point-like excitations: A gauge fluxes (the paper calls one
\texttt{m\_A}). They are the dynamical versions of the endpoints of a
defect string.

\subsection{Why the braiding is
non-Abelian}\label{why-the-braiding-is-non-abelian}

Taking an \texttt{m\_B} anyon once around an A gauge flux is equivalent
to crossing its defect string once. Hence

\[
m_B \quad\xrightarrow{\text{braid around }m_A}\quad m_B e_C.
\]

For Abelian anyons, braiding returns the same anyon type times a complex
phase, such as \texttt{-1}. Here it acts on the two-element orbit

\[
\{m_B,\;m_Be_C\}.
\]

Braiding therefore acts as a genuine transformation on an internal
space, not as one number. After gauging, the members of this orbit are
internal states of a new gauge-invariant particle. This is the intuitive
origin of its non-Abelian character (and, in the full theory, quantum
dimension greater than one).

The new A flux is special not merely because it is an extra anyon:
braiding it with B/C anyons \emph{enacts the old symmetry} on them. This
is the ``intertwining'' described in the paper.

\subsection{The torus thought
experiment}\label{the-torus-thought-experiment}

The authors next create two pairs of nontrivially braiding anyons and
wind them around the two independent cycles of a torus. The windings are
linked. In a non-Abelian theory, their order and fusion channel matter;
fusing the second pair can leave a nontrivial leftover charge at that
location. This is the ``single unpaired anyon on a torus'' image in
\texttt{single\_anyon.pdf}.

This does not violate charge conservation. The topology and the
non-Abelian fusion space store the compensating information. The lesson
is simply that the history cannot be reduced to the Abelian rule ``each
pair independently fuses back to vacuum.'' Later, the authors seek an
analogue of this phenomenon for qLDPC codes without an ordinary spatial
torus.

\subsection{\texorpdfstring{Why the result is
\(\mathcal D(D_4)\)}{Why the result is \textbackslash mathcal D(D\_4)}}\label{why-the-result-is-mathcal-dd_4}

The initial two toric codes can be viewed as an Abelian
\texttt{Z\_2\ x\ Z\_2} gauge theory. Gauging the extra global
\texttt{Z\_2} symmetry \texttt{g} produces the quantum double
\texttt{D(D\_4)}; the paper uses \texttt{D\_4} for the eight-element
dihedral group. The group language records that the old gauge structure
and the new symmetry do not remain three independent commuting
\texttt{Z\_2}s: the new operation acts nontrivially on the old one.

Two useful descriptions of the same group appear across the local
papers:

\[
D_4 \cong (\mathbb Z_2\times\mathbb Z_2)\rtimes\mathbb Z_2,
\qquad
1\to\mathbb Z_2\to D_4\to\mathbb Z_2^2\to1.
\]

The first matches this paper's ``gauge \texttt{Z\_2\ x\ Z\_2}, then
add/gauge a \texttt{Z\_2} symmetry'' viewpoint. The second is the
central-extension description made explicit in the reference
\emph{Hierarchy of topological order from finite-depth unitaries,
measurement, and feedforward}. They describe the same \texttt{D\_4}, not
two different phases.

\subsection{What to retain before Section
2.2}\label{what-to-retain-before-section-2.2}

\begin{enumerate}
\def\labelenumi{\arabic{enumi}.}
\tightlist
\item
  Two toric codes are the Abelian starting point.
\item
  Controlled-Z maps \texttt{m\_B} to \texttt{m\_B\ e\_C} and
  \texttt{m\_C} to \texttt{m\_C\ e\_B}.
\item
  A boundary of a \emph{partial} symmetry application is a defect line
  across which anyon type changes.
\item
  Gauging promotes that defect line to a physical gauge-field string
  with new flux excitations at open ends.
\item
  Braiding around one such flux changes an anyon's internal state. That,
  rather than a mere braiding phase, is the promised non-Abelian
  behavior.
\end{enumerate}

Section 2.2 implements these five statements with extra A-sublattice
qubits, cluster states, a symmetry-enrichment circuit, and
projection/measurement.

\subsection{Sources consulted locally}\label{sources-consulted-locally}

\begin{itemize}
\tightlist
\item
  Target subsection and setup:
  \href{../arXiv-2602.12228v1/main.tex}{\texttt{arXiv-2602.12228v1/main.tex}},
  Section 2 and Section 2.1 (from about line 421).
\item
  Cited background source:
  \href{../arXiv-2209.06202v2/main.tex}{\texttt{arXiv-2209.06202v2/main.tex}},
  especially Section 3.2 on one-shot \texttt{D(D\_4)} preparation and
  its \texttt{D\_4} extension data.
\end{itemize}

\end{document}
